\hypertarget{summary}{%
\section{Summary}\label{summary}}

\texttt{pyCC.id} is a Python package for data-driven system
identification that provides a flexible ``grey-box'' framework for
discovering interpretable, nonlinear governing equations from
time-series data. It is designed to bridge the gap between opaque
black-box models (e.g., Neural ODEs) and restrictive, library-dependent
methods (e.g., SINDy \autocite{Brunton2016}).

The core philosophy of \texttt{pyCC.id} is to decompose a complex
dynamical system, \(d\mathbf{x}/dt = \mathbf{F}(\mathbf{x}, t)\), into a
user-defined model structure, \(\mathbf{G}\). This structure explicitly
separates measured inputs (like system state \(\mathbf{x}\) and external
forces) from the unknown model components. These components are: (1) a
set of unknown one-dimensional \textbf{Characteristic Curves}
(\(\{\mathbf{f}\}\)), which capture the system's underlying
nonlinearities (such as stiffness or damping), and (2) a set of unknown
scalar parameters (\(\mathbf{a}\)).

The package offers a flexible, multi-backend approach to identify these
unknown components. Users can model the characteristic curves
\(\{\mathbf{f}\}\) using powerful non-biased approximators like neural
networks (NN-CC, via \texttt{PyTorch} \autocite{Paszke2019pytorch}),
simple polynomial basis functions (Poly-CC), or discover analytical
expressions directly using symbolic regression (SymbR-CC, via
\texttt{PySR} \autocite{Cranmer2023PySR}). This framework allows the
user to provide physical prior knowledge as constraints (e.g.,
symmetries, conservation laws), guiding the discovery process toward
physically consistent models for experts in engineering, physics, and
biology.

\hypertarget{statement-of-need}{%
\section{Statement of Need}\label{statement-of-need}}

Researchers in science and engineering often face a trade-off in system
identification. \textbf{Black-box models}, such as standard Neural ODEs,
can achieve high predictive accuracy but are opaque, offering little
physical insight. Conversely, \textbf{interpretable ``white-box''
methods}, like SINDy \autocite{Brunton2016} or pure Symbolic Regression
\autocite{Cranmer2023PySR}, aim to find simple analytical equations.
However, their success often depends on the true dynamics being sparsely
represented in a pre-defined library of candidate functions or, in the
case of SR, can be computationally challenging and highly sensitive to
hyperparameter tuning when applied to full, complex systems.

\texttt{pyCC.id} is built to fill this ``grey-box'' gap, targeting
domain experts who possess partial physical knowledge of their system
(i.e., the model structure \(\mathbf{G}\)) but need to discover the
specific functional forms of its components (the characteristic curves
\(\{\mathbf{f}\}\)). It addresses the need for a tool that can leverage
powerful, non-biased function approximators (like neural networks)
within a physically constrained, interpretable framework, rather than
forcing an all-or-nothing choice between accuracy and interpretability.

\hypertarget{comparison-to-other-packages}{%
\section{Comparison to other
packages}\label{comparison-to-other-packages}}

\texttt{pyCC.id} integrates concepts from several existing tools but
applies them within a unique, structured framework:

\begin{itemize}
\item
  \textbf{SINDy (e.g., \texttt{pysindy})}: SINDy \autocite{Brunton2016}
  excels when the true dynamics are a sparse combination of terms from a
  user-provided candidate library. \texttt{pyCC.id} differs by not
  relying on a pre-defined library for the full dynamics. Instead, its
  NN-CC backend learns the shape of unknown 1D functions, which can then
  be analyzed, offering flexibility when the underlying functions (e.g.,
  complex friction or stiffness) are not easily represented by simple
  library terms.
\item
  \textbf{Symbolic Regression (e.g., \texttt{PySR})}: While
  \texttt{pyCC.id} uses \texttt{PySR} \autocite{Cranmer2023PySR} as a
  backend (SymbR-CC) and a post-processing tool, its application is
  distinct. Rather than applying SR directly to the full,
  high-dimensional derivative data (a computationally difficult task),
  \texttt{pyCC.id}'s primary workflow first isolates the unknown
  components as simple 1D functions (using NN-CC) and \emph{then}
  applies SR to these much simpler, cleaner 1D curves to find their
  analytical forms.
\item
  \textbf{Black-Box Neural ODEs}: Standard Neural ODE packages learn the
  entire derivative function \(\mathbf{F}\) as a single, monolithic
  neural network. \texttt{pyCC.id} employs a ``grey-box'' approach,
  using \texttt{PyTorch} \autocite{Paszke2019pytorch} to model only the
  specific, interpretable 1D components \(\{\mathbf{f}\}\) within a
  user-defined physical structure \(\mathbf{G}\), making the resulting
  model inherently interpretable.
\end{itemize}

\hypertarget{formalism}{%
\section{Formalism}\label{formalism}}

For many physical systems, the dynamics are described by a set of
first-order ordinary differential equations (ODEs):

\[
\frac{d\mathbf{x}}{dt} = \mathbf{F}(\mathbf{x}, t)
\]

Here, \(\mathbf{x}(t)\) is the system state vector, but the function
\(\mathbf{F}\) is often a complex ``black box,'' making it difficult to
gain physical insight. The core philosophy of \texttt{pyCC.id} is to
decompose this complex function \(\mathbf{F}\) into a combination of
simpler, interpretable building blocks. We express this decomposition
as:

\[
\frac{d\mathbf{x}}{dt} = \mathbf{G}(\mathbf{x}, \mathbf{F}_{ext}(t); \{\mathbf{f}\}, \mathbf{a})
\]

Where \(\mathbf{x}\) and \(\mathbf{F}_{ext}(t)\) are the measured inputs
(system state and external forces). The goal is to find the model
components: \(\{\mathbf{f}\}\), a set of unknown 1D
\textbf{Characteristic Curves} (e.g., nonlinear stiffness or damping
functions), and \(\mathbf{a}\), a vector of unknown scalar parameters.
The user proposes the \textbf{model structure} \(\mathbf{G}\), which
defines how these components are combined. \texttt{pyCC.id} is a Python
package that discovers the optimal functions \(\{\mathbf{f}\}\) and
parameters \(\mathbf{a}\) that best fit the observed data.

\hypertarget{features}{%
\section{Features}\label{features}}

\texttt{pyCC.id} is designed to be a flexible and high-performance tool
for researchers and practitioners. Its key features include:

\begin{itemize}
\tightlist
\item
  \textbf{Interpretable Models}: Decomposes complex, high-dimensional
  dynamics into a set of simple, 1D characteristic curves that often
  have direct physical meaning (e.g., stiffness, damping).
\item
  \textbf{Flexible Function Parametrization}: Supports multiple backends
  for modeling the characteristic curves, allowing users to choose the
  right tool for their problem:

  \begin{itemize}
  \tightlist
  \item
    \textbf{Neural Networks (NN-CC)}: Uses \texttt{PyTorch}
    \autocite{Paszke2019pytorch} for high-performance, non-biased
    function approximation.
  \item
    \textbf{Polynomials (Poly-CC)}: Provides a simple baseline using
    polynomial basis functions.
  \item
    \textbf{Symbolic Regression (SymbR-CC)}: Uses \texttt{PySR}
    \autocite{Cranmer2023PySR} to discover analytical expressions
    directly.
  \end{itemize}
\item
  \textbf{Physics-Informed Discovery}: Allows users to inject domain
  knowledge as constraints during training (e.g.,
  \texttt{\textquotesingle{}f1\ odd\textquotesingle{}},
  \texttt{\textquotesingle{}f2(0)=0\textquotesingle{}}) or by defining
  conserved quantities in the loss function. This leads to more robust
  and physically consistent models.
\item
  \textbf{Hardware Acceleration}: Natively supports multicore CPUs and
  GPUs from both NVIDIA (\texttt{cuda}) and Intel (\texttt{xpu} via
  \texttt{intel-extension-for-pytorch}) for accelerating neural network
  training.
\item
  \textbf{Built-in Simulator}: Includes a versatile ODE simulator
  (\texttt{pycc.simulate}) compatible with all identified model types
  for validation and analysis.
\item
  \textbf{Comprehensive Documentation}: Provides a Google Colab notebook
  for a quick start, as well as a full gallery of tutorials and examples
  in the repository.
\end{itemize}

\hypertarget{acknowledgements}{%
\section{Acknowledgements}\label{acknowledgements}}

This work was partially supported by CONICET (Consejo Nacional de
Investigaciones Científicas y Técnicas, Argentina) under Project PIP
1679.

\hypertarget{mentions-of-scholarly-publications}{%
\section{Mentions of scholarly
publications}\label{mentions-of-scholarly-publications}}

The \texttt{pyCC.id} package serves as a generalized implementation of
the hybrid identification methodology introduced for first-order systems
in \autocite{Gonzalez2023} and \autocite{Gonzalez2024}, and later
extended to second-order systems using neural networks in
\autocite{Gonzalez2025nody} and \autocite{Gonzalez2025arxiv}.

\hypertarget{key-references}{%
\section{Key References}\label{key-references}}

The software archive is available at:
\url{https://github.com/FedejGon/pyCC.id}

\hypertarget{references}{%
\section{References}\label{references}}
